% ============================================================
% M3 S3 练习件·W2 臂D —— 课时16 2.7.2抛物线性质（照抄课时01 母版体例）
% 生成：M3 成卷轮 S3 W2 臂D｜2026-09-14｜编译 cwd＝本目录（中文路径禁 \input 绝对路径）
%   xelatex main-true.tex ×2 → 含详解印本；xelatex main-false.tex ×2 → 纯题版（单源双本）
% 输入链（全只读）：题面库 课时16-2.7.2抛物线性质.md（题面逐字装配）＋ -答案侧.md（值逐字回嵌）
%   ＋ 导学件16 main.tex 同键 ansblock（TeX 化详解底稿）；toolchain＝qp-m3.sty（钉后版
%   md5 7c3930362be8a0a2bdf21bbf8ac16573，本地挂载副本，破损即停）。
% 键账：件内 ansblock 键集 ≡ 题面库片练键集 ≡ manifest 练键序 01~16（16 键，零缺漏零多余）；
%   拓区隔离：2章-拓/测/滚 键不入本件（S4 拓展册收）；导学键 G1~G5 属导学件域，亦不入。
% 卷式例外案B：练习件不属卷式——答案块物理落点恒为题后 ansblock，禁卷末附卷制；
%   全线括线模（导言 \ansblockgrayfalse 一次置定，件内不切换；灰底盒不可跨栏断，练习件禁用）。
% 题序＝manifest 键序 01~16 零重排；印面题号 1~16 连号（\tihao 号＝\ansitem 号同键）；
%   三档切位 1~7／8~14／15~16（照库序切位，非难度重排）。
% 槽配实录：单选3（题9/10/11）＋多选1（题1，≤4 合规）＋填空0＋解答12（题2~8、12~16）；
%   照题面侧头标（批D §四 满足度表同口径）。
% 选项槽位：默认四连排（0.25\linewidth−0.25\lxhang），超宽降档梯 四连排→两连排
%   （0.5\linewidth−0.5\lxhang−0.25em）→单列 \lxopt；禁缩字号/负 kern/删标点凑宽。
%   本片实配：题1/题10 单列（题1 导学件16 同款；题10 C/D 槽 34.3mm 超 strict 预警 31.8mm 降档）、
%   题9/11 四连排（降档实测见值台账＋回执）。
% 解答书写区：题2~8、12~16 十二题 \liubai（带内 16~40mm）：题3 四小题 32mm、题15 两问 24mm、
%   余默认 16mm；纯题档吞 ansblock 不吞 \liubai。
% tieside：难度×知识点照同章导学件16 同题标签（逐题登记见值台账 难度词派生标注）。
% 印面清洗（导学件底稿→练习件印面）：剔 \ding 记号×4（题2/5/7 检验句文句保留）；
%   题16「（到准线距离，焦半径公式见例 36）」→「（即到准线的距离）」（拓册引用改文句）；
%   题面侧【注】行不入印面（本片练键零注）。
% ============================================================
\documentclass[fontset=none]{ctexart}
\usepackage{qp-m3}
% —— 印面逐字符号路由补钉（承导学件母版；− 走 TNR、▱ NSC 子块；禁改 sty 保 md5） ——
\xeCJKDeclareCharClass{Default}{"2212}%
\setCJKfamilyfont{nscblk}[Path=C:/提示词/工作区/字替对照-0909/variantF/fonts/,BoldFont=NSC-w600.ttf]{NSC-w500.ttf}%
\xeCJKDeclareSubCJKBlock{nscblk}{"25B1}%
\newcommand{\pxparallelogram}{{\CJKfamily{nscblk}▱}}%
% —— 断行微调（承导学件母版实测档，三零门承重；\emergencystretch 不另设） ——
\xeCJKsetup{CJKglue={\hskip 0.10em plus 0.08em minus 0.05em}}%
% —— 练习件全线括线模（S3 波次方案 §四.3：导言一次置定、件内不切换） ——
\ansblockgrayfalse
% —— 页脚件名（qp-headfoot 复用入口） ——
\renewcommand{\qpzhangming}{第二章\quad 平面解析几何}%
\renewcommand{\qpjianming}{练习件}%

\begin{document}
%装配轮S6三本跨本连续页码（G7方案B）setcounter 注入位
\setcounter{page}{110}

\keshi{第16课时\quad 2.7.2抛物线性质}
\begin{multicols}{2}
\emergencystretch=1em

% ============================================================
% 基础巩固（题1~7＝16-01~16-07：多选1＋解答6；题后紧跟答案制，全线括线 ansblock）
% ============================================================
\dangbadge{基}{础}{巩}{固}

\tihao{1}\duoxuan\hspace{0.6em}\zu{性质多选辨析×多选} \tieside{简单(知识点一)}关于抛物线\(y^{2}=-2x\)，下列说法正确的是\nobreak(\kongwei)

\lxopt{A．开口向左}

\lxopt{B．焦点坐标为\((-1,0)\)}

\lxopt{C．准线为 \(x=1\)}

\lxopt{D．对称轴为 \(x\) 轴}

\begin{ansblock}[2章-练-课时16-01]
% ans:2章-练-课时16-01
\ansitem{1}{AD}
\ansnote{详解}{对照标准形 \(y^{2}=-2px\)（\(p>0\)）：\(2p=2\)、\(p=1\)，一次项系数为负，开口向左，A 正确；焦点为 \((-p/2,0)=(-1/2,0)\)，B 所给 \((-1,0)\) 错（把 \(p\) 误作 \(p/2\)）；准线为 \(x=p/2=1/2\)，C 所给 \(x=1\) 错；对称轴为 \(x\) 轴，D 正确．故选：AD．}
\end{ansblock}

\tihao{2}\zu{焦半径与定义互化×解答} \tieside{简单(知识点二)}已知抛物线 \(y=4x^{2}\) 上的一点 \(M\) 的纵坐标为 1，求点 \(M\) 到焦点的距离．
\liubai[16mm]{解答书写区}

\begin{ansblock}[2章-练-课时16-02]
% ans:2章-练-课时16-02
\ansitem{2}{17/16}
\ansnote{详解}{化标准形：\(y=4x^{2}\) 即 \(x^{2}=(1/4)y\)，故 \(2p=1/4\)、\(p=1/8\)，焦点 \(F(0,1/16)\)，准线 \(y=-1/16\)．由抛物线定义，点 \(M\) 到焦点的距离等于它到准线的距离，\(|MF|=y_{M}+p/2=1+1/16=17/16\)．（回代校验：\(x_{M}^{2}=1/4\)、\(M(\pm 1/2,1)\)，\(|MF|^{2}=x_{M}^{2}+(y_{M}-1/16)^{2}=1/4+225/256=289/256\)，\(|MF|=17/16\)．）}
\end{ansblock}

\tihao{3}\zu{焦点与准线×解答} \tieside{简单(知识点一)}求下列方程表示的抛物线的焦点坐标和准线方程：
(1) \(y^{2}-6x=0\)；　(2) \(x^{2}+10y=0\)；
(3) \(y=(3/8)x^{2}\)；　(4) \(ax+y^{2}=0\)（\(a\neq 0\)）．
\liubai[32mm]{解答书写区}

\begin{ansblock}[2章-练-课时16-03]
% ans:2章-练-课时16-03
\ansitem{3}{(1) 焦点(3/2,0)，准线 x=−3/2；(2) 焦点(0,−5/2)，准线 y=5/2；(3) 焦点(0,2/3)，准线 y=−2/3；(4) 焦点(−a/4,0)，准线 x=a/4}
\ansnote{详解}{(1) \(y^{2}=6x\)：\(2p=6\)、\(p=3\)，开口向右，焦点 \((p/2,0)=(3/2,0)\)，准线 \(x=-3/2\)．(2) \(x^{2}=-10y\)：\(2p=10\)、\(p=5\)，开口向下，焦点 \((0,-5/2)\)，准线 \(y=5/2\)．(3) \(x^{2}=(8/3)y\)：\(2p=8/3\)、\(p=4/3\)，开口向上，焦点 \((0,2/3)\)，准线 \(y=-2/3\)．(4) \(y^{2}=-ax\)（\(a\neq 0\)）：一次项系数为 \(-a\)，无论 \(a>0\)（开口向左）或 \(a<0\)（开口向右），焦点均为 \((-a/4,0)\)、准线均为 \(x=a/4\)（\(a>0\) 时焦点在 \(x\) 轴负半轴，\(a<0\) 时在正半轴，公式统一）．}
\end{ansblock}

\tihao{4}\zu{焦点与准线×解答} \tieside{简单(知识点一)}求抛物线 \(y^{2}=8x\) 的焦点到直线 \(x-\sqrt{3}y=0\) 的距离．
\liubai[16mm]{解答书写区}

\begin{ansblock}[2章-练-课时16-04]
% ans:2章-练-课时16-04
\ansitem{4}{1}
\ansnote{详解}{\(2p=8\)、\(p=4\)，焦点 \(F(2,0)\)．由点到直线的距离公式，\(d=|2-\sqrt{3}\times 0|/\sqrt{1^{2}+(\sqrt{3})^{2}}=2/\sqrt{4}=1\)．}
\end{ansblock}

\tihao{5}\zu{焦半径与定义互化×解答} \tieside{简单(知识点二)}设抛物线 \(y^{2}=8x\) 上一点 \(P\) 到 \(y\) 轴的距离是 4，求点 \(P\) 到该抛物线焦点的距离．
\liubai[16mm]{解答书写区}

\begin{ansblock}[2章-练-课时16-05]
% ans:2章-练-课时16-05
\ansitem{5}{6}
\ansnote{详解}{\(y^{2}=8x\) 中 \(2p=8\)、\(p=4\)，准线 \(x=-2\)．点 \(P\) 到 \(y\) 轴距离为 4，而抛物线上点横坐标 \(x\geqslant 0\)，故 \(x_{P}=4\)．由定义，\(P\) 到焦点的距离等于 \(P\) 到准线的距离：\(|PF|=x_{P}+p/2=4+2=6\)．（校验：\(P(4,\pm 4\sqrt{2})\)，\(|PF|=\sqrt{(4-2)^{2}+32}=\sqrt{36}=6\)．）}
\end{ansblock}

\tihao{6}\zu{焦半径反求坐标×解答} \tieside{简单(知识点二)}若抛物线 \(y^{2}=2px\)（\(p>0\)）上一点 \(M\) 到焦点 \(F\) 的距离 \(|MF|=2p\)，求点 \(M\) 的坐标．
\liubai[16mm]{解答书写区}

\begin{ansblock}[2章-练-课时16-06]
% ans:2章-练-课时16-06
\ansitem{6}{M(3p/2, ±√3 p)}
\ansnote{详解}{设 \(M(x_{0},y_{0})\)，\(x_{0}\geqslant 0\)．由定义，\(|MF|\) 等于 \(M\) 到准线 \(x=-p/2\) 的距离，即 \(x_{0}+p/2=2p\)，得 \(x_{0}=3p/2\)．代入方程：\(y_{0}^{2}=2p\cdot(3p/2)=3p^{2}\)，\(y_{0}=\pm\sqrt{3}p\)．故 \(M(3p/2,\sqrt{3}p)\) 或 \(M(3p/2,-\sqrt{3}p)\)．}
\end{ansblock}

\tihao{7}\zu{依点求方程×解答} \tieside{简单(知识点一)}已知抛物线的顶点在坐标原点，对称轴是坐标轴，且过点 \(P(-2,2\sqrt{2})\)，求抛物线的标准方程．
\liubai[16mm]{解答书写区}

\begin{ansblock}[2章-练-课时16-07]
% ans:2章-练-课时16-07
\ansitem{7}{y²=−4x 或 x²=√2 y}
\ansnote{详解}{\(P(-2,2\sqrt{2})\) 在第二象限，开口只可能向左或向上．情形一，对称轴为 \(x\) 轴：设 \(y^{2}=-2px\)（\(p>0\)），代入得 \((2\sqrt{2})^{2}=-2p\cdot(-2)=4p\)、\(p=2\)，即 \(y^{2}=-4x\)．情形二，对称轴为 \(y\) 轴：设 \(x^{2}=2py\)（\(p>0\)），代入得 \(4=2p\cdot 2\sqrt{2}\)、\(2p=\sqrt{2}\)，即 \(x^{2}=\sqrt{2}y\)．（回验：\(8=-4\times(-2)\)；\(4=\sqrt{2}\times 2\sqrt{2}\)．）}
\end{ansblock}

% ============================================================
% 综合提升（题8~14＝16-08~16-14：解答4＋单选3，书写区 \liubai＋题后括线 ansblock）
% ============================================================
\dangbadge{综}{合}{提}{升}

\tihao{8}\zu{跨曲线联结求参数×解答} \tieside{简单(知识点一)}若抛物线 \(y^{2}=2px\)（\(p>0\)）的焦点与椭圆 \(x^{2}/6+y^{2}/2=1\) 的右焦点重合，求 \(p\) 的值．
\liubai[16mm]{解答书写区}

\begin{ansblock}[2章-练-课时16-08]
% ans:2章-练-课时16-08
\ansitem{8}{p=4}
\ansnote{详解}{椭圆中 \(c^{2}=6-2=4\)，右焦点 \((2,0)\)．抛物线焦点 \((p/2,0)\)，由 \(p/2=2\) 得 \(p=4\)．}
\end{ansblock}

\tihao{9}\zu{焦点弦×单选} \tieside{简单(知识点三)}已知抛物线 \(C\)：\(y^{2}=x\) 的焦点为 \(F\)，过点 \(F\) 且斜率为 \(\sqrt{3}\) 的直线交抛物线 \(C\) 于 \(A\)、\(B\) 两点，则 \(|AF|\cdot|BF|\) 等于\nobreak(\kongwei)
\bindopt {\fontsize{10.5pt}{\qpTJgLeadLiti}\selectfont \makebox[\dimexpr0.25\linewidth-0.25\lxhang\relax][l]{A．\(1/3\)}\makebox[\dimexpr0.25\linewidth-0.25\lxhang\relax][l]{B．\(4/3\)}\makebox[\dimexpr0.25\linewidth-0.25\lxhang\relax][l]{C．\(1\)}\makebox[\dimexpr0.25\linewidth-0.25\lxhang\relax][l]{D．\(4\)}\par}

\begin{ansblock}[2章-练-课时16-09]
% ans:2章-练-课时16-09
\ansitem{9}{A（1/3）}
\ansnote{详解}{\(y^{2}=x\) 中 \(2p=1\)，焦点 \(F(1/4,0)\)，准线 \(x=-1/4\)．直线 \(l\)：\(y=\sqrt{3}(x-1/4)\)．代入 \(y^{2}=x\) 得 \(3(x-1/4)^{2}=x\)，整理为 \(48x^{2}-40x+3=0\)，故 \(x_{1}+x_{2}=40/48=5/6\)，\(x_{1}x_{2}=3/48=1/16\)．由定义，\(|AF|=x_{1}+1/4\)、\(|BF|=x_{2}+1/4\)，\(|AF|\cdot|BF|=x_{1}x_{2}+(1/4)(x_{1}+x_{2})+1/16=1/16+5/24+1/16=16/48=1/3\)．故选 A．}
\end{ansblock}

\tihao{10}\zu{通径与标准方程×单选} \tieside{简单(知识点一)}以 \(x\) 轴为对称轴，抛物线通径的长为 8，顶点在坐标原点的抛物线的方程是\nobreak(\kongwei)
\lxopt{A．\(y^{2}=8x\)}

\lxopt{B．\(y^{2}=-8x\)}

\lxopt{C．\(y^{2}=8x\) 或 \(y^{2}=-8x\)}

\lxopt{D．\(x^{2}=8y\) 或 \(x^{2}=-8y\)}

\begin{ansblock}[2章-练-课时16-10]
% ans:2章-练-课时16-10
\ansitem{10}{C}
\ansnote{详解}{顶点在原点、对称轴为 \(x\) 轴，方程形如 \(y^{2}=\pm 2px\)（\(p>0\)）．通径（过焦点且垂直于对称轴的弦）长为 \(2p\)，由 \(2p=8\) 得 \(p=4\)，故方程为 \(y^{2}=8x\)（焦点在 \(x\) 轴正半轴）或 \(y^{2}=-8x\)（焦点在负半轴），两向各一，选 C．（排除 D：其对称轴为 \(y\) 轴．）}
\end{ansblock}

\tihao{11}\zu{弦长与准线距离×单选} \tieside{中档(知识点三)}已知 \(F\) 为抛物线 \(C\)：\(y^{2}=4x\) 的焦点，过点 \(F\) 的直线 \(l\) 交抛物线 \(C\) 于 \(A\)，\(B\) 两点．若 \(|AB|=8\)，则线段 \(AB\) 的中点 \(M\) 到直线 \(x+1=0\) 的距离为\nobreak(\kongwei)
\bindopt {\fontsize{10.5pt}{\qpTJgLeadLiti}\selectfont \makebox[\dimexpr0.25\linewidth-0.25\lxhang\relax][l]{A．\(2\)；}\makebox[\dimexpr0.25\linewidth-0.25\lxhang\relax][l]{B．\(4\)；}\makebox[\dimexpr0.25\linewidth-0.25\lxhang\relax][l]{C．\(8\)；}\makebox[\dimexpr0.25\linewidth-0.25\lxhang\relax][l]{D．\(16\)}\par}

\begin{ansblock}[2章-练-课时16-11]
% ans:2章-练-课时16-11
\ansitem{11}{B}
\ansnote{详解}{\(y^{2}=4x\) 的焦点 \(F(1,0)\)，准线 \(x=-1\)，即直线 \(x+1=0\)．分别过 \(A\)、\(B\)、\(M\) 作准线的垂线，垂足为 \(C\)、\(D\)、\(N\)．由定义，\(|AC|=|AF|\)、\(|BD|=|BF|\)，故 \(|AC|+|BD|=|AF|+|BF|=|AB|=8\)．又 \(AC\parallel MN\parallel BD\) 且 \(M\) 为 \(AB\) 中点，\(MN\) 为直角梯形 \(ABDC\) 的中位线，\(|MN|=(|AC|+|BD|)/2=4\)，即点 \(M\) 到准线 \(x=-1\) 的距离为 4．故选：B．}
\end{ansblock}

\tihao{12}\zu{定义转化求最值×解答} \tieside{中档(知识点二)}已知抛物线的方程为 \(x^{2}=8y\)，\(F\) 是其焦点，点 \(A(-2,4)\) 在抛物线的内部，在此抛物线上求一点 \(P\)，使 \(|PF|+|PA|\) 的值最小．
\liubai[16mm]{解答书写区}

\begin{ansblock}[2章-练-课时16-12]
% ans:2章-练-课时16-12
\ansitem{12}{P(−2,1/2)（最小值 6）}
\ansnote{详解}{\(x^{2}=8y\) 的准线为 \(l\)：\(y=-2\)．设 \(P\) 在 \(l\) 上的射影为 \(Q\)，过 \(A\) 作 \(AB\perp l\) 于 \(B\)，则 \(|AB|=y_{A}-(-2)=4+2=6\)．由定义 \(|PF|=|PQ|\)，故 \(|PF|+|PA|=|PQ|+|PA|\geqslant|AQ|\geqslant|AB|=6\)：第一个不等号当 \(A\)、\(P\)、\(Q\) 共线（\(P\) 与 \(A\) 同取横坐标 \(-2\)）时取等，第二个不等号当 \(Q\) 与 \(B\) 重合时取等，两条件同时满足．把 \(x=-2\) 代入 \(x^{2}=8y\) 得 \(y=1/2\)，故所求点 \(P(-2,1/2)\)，此时 \((|PF|+|PA|)_{\min}=|AB|=y_{A}+2=6\)．}
\end{ansblock}

\tihao{13}\zu{曲线上动点到定点最值×解答} \tieside{中档(知识点二)}已知抛物线 \(y^{2}=16x\) 和点 \(A(4,0)\)，点 \(M\) 在此抛物线上运动，求点 \(M\) 到点 \(A\) 的距离的最小值，并指出取得最小值时点 \(M\) 的坐标．
\liubai[16mm]{解答书写区}

\begin{ansblock}[2章-练-课时16-13]
% ans:2章-练-课时16-13
\ansitem{13}{|MA| 的最小值为 4，此时 M(0,0)}
\ansnote{详解}{设 \(M(x,y)\)，则 \(x\geqslant 0\)，\(y^{2}=16x\)．\par\noindent \(|MA|^{2}=(x-4)^{2}+y^{2}=(x-4)^{2}+16x=x^{2}+8x+16=(x+4)^{2}\)．\par\noindent 由 \(x\geqslant 0\) 得 \((x+4)^{2}\geqslant 16\)，当且仅当 \(x=0\) 时等号成立，此时 \(y=0\)，\(|MA|_{\min}=4\)，\(M(0,0)\)．（注：\(2p=16\)、\(p/2=4\)，\(A(4,0)\) 恰为该抛物线的焦点，\(|MA|=|MF|=x+4\geqslant 4\) 与配方路径互证．）}
\end{ansblock}

\tihao{14}\zu{正三角形与抛物线×解答} \tieside{中档(知识点三)}正三角形的一个顶点位于坐标原点，另外两个顶点在抛物线 \(y^{2}=2px\)（\(p>0\)）上，求这个正三角形的边长．
\liubai[16mm]{解答书写区}

\begin{ansblock}[2章-练-课时16-14]
% ans:2章-练-课时16-14
\ansitem{14}{4√3 p}
\ansnote{详解}{设正三角形 \(OAB\) 的顶点 \(A(x_{1},y_{1})\)、\(B(x_{2},y_{2})\) 在抛物线上，则 \(y_{1}^{2}=2px_{1}\)、\(y_{2}^{2}=2px_{2}\)．由 \(|OA|=|OB|\) 得 \(x_{1}^{2}+y_{1}^{2}=x_{2}^{2}+y_{2}^{2}\)，即 \((x_{1}^{2}-x_{2}^{2})+2p(x_{1}-x_{2})=0\)，\((x_{1}-x_{2})(x_{1}+x_{2}+2p)=0\)．因 \(x_{1}\geqslant 0\)、\(x_{2}\geqslant 0\) 且 \(2p>0\)（若 \(x_{1}\)、\(x_{2}\) 之一为 0 则三角形退化），括号第二因子为正，故 \(x_{1}=x_{2}\)，从而 \(y_{2}=-y_{1}\neq 0\)，线段 \(AB\) 关于 \(x\) 轴对称．\(x\) 轴垂直平分 \(AB\) 且平分 \(\angle AOB\)，\(\angle AOx=30^{\circ}\)，故 \(y_{1}/x_{1}=\tan 30^{\circ}=\sqrt{3}/3\)、\(x_{1}=\sqrt{3}y_{1}\)．代入 \(y_{1}^{2}=2p\cdot\sqrt{3}y_{1}\) 得 \(y_{1}=2\sqrt{3}p\)，边长 \(|AB|=2y_{1}=4\sqrt{3}p\)．}
\end{ansblock}

% ============================================================
% 思维探索（题15~16＝16-15~16-16：解答2 压轴对，居末档照库序）
% ============================================================
\dangbadge{思}{维}{探}{索}

\tihao{15}\zu{压轴：焦点弦定值×解答} \tieside{难(知识点三)}已知 \(O\) 为坐标原点，过抛物线 \(y^{2}=2px\)（\(p>0\)）的焦点 \(F\) 作一条倾斜角为 \(\pi/4\) 的直线，与抛物线相交于 \(A\)，\(B\) 两点．
(1) 用 \(p\) 表示 \(A\)，\(B\) 之间的距离；
(2) 证明：\(\cos\angle AOB\) 的大小是与 \(p\) 无关的定值，并求出这个值．
\liubai[24mm]{解答书写区}

\begin{ansblock}[2章-练-课时16-15]
% ans:2章-练-课时16-15
\ansitem{15}{(1) |AB|=4p；(2) 定值 cos∠AOB=−3√41/41（与 p 无关）}
\ansnote{详解}{(1) 过 \(F(p/2,0)\)、倾斜角 \(\pi/4\) 的直线为 \(y=x-p/2\)．代入 \(y^{2}=2px\)：\((x-p/2)^{2}=2px\)，即 \(x^{2}-3px+p^{2}/4=0\)（\(\Delta=9p^{2}-p^{2}=8p^{2}>0\)）．设 \(A(x_{A},y_{A})\)、\(B(x_{B},y_{B})\)，则 \(x_{A}+x_{B}=3p\)，\(x_{A}x_{B}=p^{2}/4\)．由定义（焦半径），\par\noindent \(|AB|=|AF|+|BF|=(x_{A}+p/2)+(x_{B}+p/2)=x_{A}+x_{B}+p=4p\)．\par\noindent (2) 该弦所在直线与 \(x\) 轴交于 \(F\)，两端点 \(A\)、\(B\) 分居 \(x\) 轴两侧，纵坐标一正一负，其积为负：\par\noindent \(y_{A}y_{B}=(x_{A}-p/2)(x_{B}-p/2)=x_{A}x_{B}-(p/2)(x_{A}+x_{B})+p^{2}/4=p^{2}/4-3p^{2}/2+p^{2}/4=-p^{2}\)（与焦点弦结论 \(y_{1}y_{2}=-p^{2}\) 相符，此即符号判据）．\par\noindent 于是 \(x_{A}x_{B}+y_{A}y_{B}=p^{2}/4-p^{2}=-3p^{2}/4\)．又 \(|OA|^{2}=x_{A}^{2}+y_{A}^{2}=x_{A}^{2}+2px_{A}=x_{A}(x_{A}+2p)\)，同理 \(|OB|^{2}=x_{B}(x_{B}+2p)\)．\par\noindent 故 \(|OA|\cdot|OB|=\sqrt{x_{A}x_{B}}\cdot\sqrt{x_{A}x_{B}+2p(x_{A}+x_{B})+4p^{2}}=(p/2)\cdot\sqrt{p^{2}/4+6p^{2}+4p^{2}}=(p/2)\cdot(\sqrt{41}/2)p=\sqrt{41}p^{2}/4\)．\par\noindent 在 \(\triangle AOB\) 中由余弦定理，\(\cos\angle AOB=(|OA|^{2}+|OB|^{2}-|AB|^{2})/(2|OA|\cdot|OB|)\)．\par\noindent 而 \(|OA|^{2}+|OB|^{2}-|AB|^{2}=(x_{A}-x_{B})^{2}+(y_{A}-y_{B})^{2}=2(x_{A}x_{B}+y_{A}y_{B})=-3p^{2}/2\)，\par\noindent 故 \(\cos\angle AOB=(-3p^{2}/4)/(\sqrt{41}p^{2}/4)=-3/\sqrt{41}\)，分母有理化得 \(-3\sqrt{41}/41\)，\(p^{2}\) 全部约去，与 \(p\) 无关．}
\end{ansblock}

\tihao{16}\zu{压轴：弦长和最值×解答} \tieside{难(知识点三)}过抛物线 \(C\)：\(x^{2}=4y\) 的焦点 \(F\) 作斜率分别为 \(k_{1}\)、\(k_{2}\) 的两条直线 \(l_{1}\)、\(l_{2}\)，其中 \(l_{1}\) 交抛物线 \(C\) 于 \(A\)、\(B\) 两点，\(l_{2}\) 交抛物线 \(C\) 于 \(D\)、\(E\) 两点．若 \(k_{1}k_{2}=-2\)，求 \(|AB|+|DE|\) 的最小值．
\liubai[16mm]{解答书写区}

\begin{ansblock}[2章-练-课时16-16]
% ans:2章-练-课时16-16
\ansitem{16}{24}
\ansnote{详解}{\(x^{2}=4y\) 的焦点 \(F(0,1)\)，准线 \(y=-1\)，抛物线上点之焦半径 \(|AF|=y_{A}+1\)（即到准线的距离）．\(l_{1}\)：\(y=k_{1}x+1\) 代入 \(x^{2}=4y\) 得 \(x^{2}-4k_{1}x-4=0\)（\(\Delta=16k_{1}^{2}+16>0\) 恒成立）．设 \(A(x_{1},y_{1})\)、\(B(x_{2},y_{2})\)，则 \(x_{1}+x_{2}=4k_{1}\)，\(y_{1}+y_{2}=k_{1}(x_{1}+x_{2})+2=4k_{1}^{2}+2\)．于是 \par\noindent \(|AB|=|AF|+|BF|=(y_{1}+1)+(y_{2}+1)=y_{1}+y_{2}+2=4k_{1}^{2}+4\)．\par\noindent 同理 \(|DE|=4k_{2}^{2}+4\)．故 \par\noindent \(|AB|+|DE|=4(k_{1}^{2}+k_{2}^{2})+8\geqslant 8|k_{1}k_{2}|+8=8\times 2+8=24\)，\par\noindent 当且仅当 \(|k_{1}|=|k_{2}|\) 且 \(k_{1}k_{2}=-2\)，即 \((k_{1},k_{2})=(\sqrt{2},-\sqrt{2})\) 或 \((-\sqrt{2},\sqrt{2})\) 时等号成立．所求最小值为 24．}
\end{ansblock}

% ---- 尾块（\tailfill 置 \end{multicols} 前；无参＝内容尾随框，每件恰 1 处） ----
\tailfill

\end{multicols}
\end{document}
